% Separate CE6190 report; official CVPR author-kit files remain unmodified.
\documentclass[10pt,twocolumn,letterpaper]{article}
\usepackage{cvpr}
%% This file contains a number of tweaks that are typically applied to the main document.
%% They are not enabled by default, but can be enabled by uncommenting the relevant lines.


%%
%% Official support for 2-column teaser figure
%%
\usepackage{cuted} %< for strip environment

%%
%% For auto-referencing of figures (see fig/teaser.tex and \label{} therein) 
%%
\usepackage{currfile} %< for \currfilebase macro

%%
%% caption and (global) spacing overrides
%%
\usepackage{caption} %< for "\captionof" commands (teaser)
% \setlength{\abovecaptionskip}{.5em} % space above caption
% \setlength{\belowcaptionskip}{1em} % space below caption

%%
%% Inline annotations; for predefined colors, refer to "dvipsnames" in the xcolor package:
%% https://tinyurl.com/overleaf-colors
%%
\newcommand{\red}[1]{{\color{red}#1}}
\newcommand{\todo}[1]{{\color{red}#1}}
\newcommand{\TODO}[1]{\textbf{\color{red}[TODO: #1]}}
%%
%% Disable for camera ready / submission by uncommenting these lines:
%%
% \renewcommand{\TODO}[1]{}
% \renewcommand{\todo}[1]{#1}

%%
%% Work harder in optimizing text layout. Typically shrinks text by 1/6 of page, enable
%% it at the very end of the writing process, when you are just above the page limit.
%%
% \usepackage{microtype}

%%
%% Fine-tune paragraph spacing.
%%
% \renewcommand{\paragraph}[1]{\vspace{.5em}\noindent\textbf{#1.}}

%%
%% Globally adjusts space between figure and caption.
%%
% \setlength{\abovecaptionskip}{.5em}

%%
%% Allows "the use of \paper to refer to the project name"
%% with automatic management of space at the end of the word.
%%
% \usepackage{xspace}
% \newcommand{\paper}{ProjectName\xspace}

%%
%% Commonly used math definitions.
%%
% \DeclareMathOperator*{\argmin}{arg\,min}
% \DeclareMathOperator*{\argmax}{arg\,max}

%%
%% Tighten underline.
%%
% \usepackage{soul}
% \setuldepth{foobar}

%%
%% Reconfigure paragraph (tweaks spacing)
%%

\usepackage[hidelinks]{hyperref}
\usepackage{tabularx}
\hypersetup{pdftitle={One-Step Text-to-Image Distillation: Comparing Standard Pipelines and Direction Alignment},
pdfauthor={Le Duc Dong; Nguyen Huu Thuat; Pham Viet Hoang; La Nguyen Gia Hy},
pdfsubject={CE6190 Assignment 2, Group 11}}
\def\paperID{}
\def\confName{CVPR}
\def\confYear{2026}
\input{generated/assignment2/group.tex}
\input{generated/assignment2/settings.tex}
\input{generated/assignment2/main_selection.tex}
\input{generated/assignment2/selected_response_metrics.tex}
\newcommand{\sg}{\operatorname{sg}}
\title{One-Step Text-to-Image Distillation:\\Comparing Standard Pipelines and Direction Alignment}
\author{Le Duc Dong \quad Nguyen Huu Thuat \quad Pham Viet Hoang \quad La Nguyen Gia Hy\\
CE6190: Recent Advance on Visual Understanding and Generation\\
Assignment 2: Visual Generation \quad Group \AssignmentGroup}
\begin{document}
\maketitle

\begin{abstract}
One-step distillation compresses a diffusion sampler into a single network
call, but its learning signal and training cost vary substantially across
methods. We compare progressive distillation, consistency distillation,
distribution matching distillation (DMD), SwiftBrush and InstaFlow using
Stable Diffusion 1.5 with rank-32 generator adapters. All five pipelines
receive 36,000 updates and are evaluated on the same 5,000 COCO-2017
validation captions at one network function evaluation. DMD with its
original perceptual regression is the strongest standard pipeline, with
FID \ADMDFID{}. We also investigate direction alignment, which adds a
constraint on how paired student outputs respond to different noise inputs.
Our independently trained model achieves FID \ASelectedFID{} after
\ASelectedUpdates{} updates, with lower IS and CLIP than the original DMD
reference. Matched native samples and a 36,000-update regression control accompany
the quantitative comparison. The original reference uses a different
initialization history and teacher bank, so this comparison measures the
complete recipes. Detailed design and ablations are provided in the appendix.
\end{abstract}

\section{Introduction}
Text-to-image diffusion models generate detailed images by repeatedly
refining noise under text conditioning. Latent diffusion reduces the cost
of each refinement by operating in a learned image representation, but a
high-quality sample can still require many network
calls~\cite{rombach2022ldm}. One-step distillation instead learns a direct
mapping from noise and a caption to an image. This makes the training
objective responsible for structure and detail that would otherwise emerge
through a sequence of denoising updates.

Several approaches address this compression differently. Progressive
distillation shortens a teacher trajectory in stages; consistency
distillation learns a common endpoint along a trajectory. DMD and
SwiftBrush use diffusion scores to improve the generated distribution,
while InstaFlow uses reflow to simplify noise--image paths. These methods
also differ in their dependence on paired targets, auxiliary networks and
teacher sampling. An inference-step comparison alone does not capture
these training costs.

We study these choices under a practical project budget: a shared SD1.5
backbone, rank-32 student adapters and 36,000 updates per standard pipeline.
Our first question is which local recipe gives the best quality at one
network function evaluation (NFE). Our second is what the comparison
reveals about supervision: in particular, how much DMD's original
perceptual regression contributes relative to latent regression.

We then examine an addition to the strongest standard pipeline. Pointwise
regression anchors each generated image to a teacher output, but it does
not explicitly align the direction of the change between two noise inputs.
Our direction-alignment objective adds this constraint while retaining
DMD's distribution-matching and perceptual terms. The completed 36,000-update result is
included as \textbf{Ours} in the main comparison; its mathematical design,
fixed training settings and independent control are in
Appendix~\ref{sec:response}.

The study contributes a measured comparison of five one-step recipes,
a matched regression-space analysis, and an evaluated direction-alignment
extension with native visualizations. We distinguish matched inference
from matched training compute and report the proposal's actual update
count. The resulting conclusions concern this small-adapter setting rather
than the published full-model rankings.

\begin{figure*}[t]
\centering
\includegraphics[width=\textwidth]{generated/assignment2/main_samples.pdf}
\caption{\textbf{One-step comparison including our selected proposal.}
Columns show the five standard pipelines, \textbf{Ours} (direction alignment),
and the SD1.5 teacher. Each row uses the same caption and initial noise.
All students use one step / one NFE; the teacher uses 64 DDIM steps / 128 NFE.
All student generators have 36,000 updates.
Dog and lighthouse prompts are preset, and teddy bears are a selected
validation example. The illustrations are separate from the full-set metrics.}
\label{fig:main}
\end{figure*}

\section{Methods}
\subsection{Common task and model}
For caption $c$ and Gaussian noise $z\sim\mathcal N(0,I)$, the student
produces latent $G_\theta(z,c)$ and image $D(G_\theta(z,c))$, where $D$ is the
frozen VAE decoder. Images have resolution $512\times512$. The base U-Net,
text encoder and VAE remain frozen; the trainable generator parameters are
SD1.5 attention LoRA adapters~\cite{hu2022lora},
$W'=W+(a/r)BA$, with rank $r=32$. This gives 6,377,472 trainable generator
parameters per method.

All students absorb teacher guidance during training and use a single
conditional branch at inference. The DDIM teacher uses classifier-free
guidance with two branches per step~\cite{song2021ddim}. We count denoiser
branches as NFE, excluding text encoding and VAE decoding. Thus one student
step costs one NFE, whereas the 64-step teacher costs 128 NFE. The common
architecture controls student capacity; each method retains its own
supervision and auxiliary components.

\subsection{Five standard distillation pipelines}
\paragraph{Progressive distillation (PD).}
PD trains one student transition to match two teacher transitions and
repeats the compression~\cite{salimans2022progressive}. Our schedule is
$64\!\rightarrow\!32\!\rightarrow\!16\!\rightarrow\!8\!\rightarrow\!4
\!\rightarrow\!2\!\rightarrow\!1$, using clipped-SNR weighting at threshold
5. The 36,000-update budget covers the whole pipeline. Local transition
targets make the learning problem explicit, but each stage depends on the
quality of the preceding distilled trajectory. Errors can therefore
propagate as the sampling budget decreases.

\paragraph{Consistency distillation.}
Consistency models map points on the same probability-flow trajectory to
a common endpoint~\cite{song2023consistency}. A teacher transition from
$x_t$ to $x_s$ supplies the representative objective
\begin{equation}
L_{\mathrm{CD}}=\mathbb E\|f_\theta(x_t,t,c)
-\sg[f_{\bar\theta}(x_s,s,c)]\|_2^2,
\end{equation}
where $\bar\theta$ is an exponential moving average (EMA) target and $\sg$
stops gradients. Our implementation uses a variance-preserving boundary
blend with $\sigma_{\mathrm{data}}=0.5$ and latent L2 supervision. These
choices adapt the method to the shared latent backbone; they differ from
the original EDM/image-perceptual recipe. The endpoint mapping is efficient
at inference, while its training depends on the boundary parameterization
and distance measure.

\paragraph{Distribution matching distillation (DMD).}
DMD combines distribution matching with paired perceptual
regression~\cite{yin2024dmd}. A frozen diffusion teacher estimates the target
score, and an auxiliary fake scorer learns from generated samples. Given
student output $x$, the scorers process the same noisy latent and predict
clean latents $\hat x_0^T$ and $\hat x_0^F$. The normalized generator
direction in our implementation is
\begin{equation}
g=\frac{\hat x_0^F-\hat x_0^T}
 {\max\{\operatorname{mean}(|x-\hat x_0^T|),\epsilon\}}.
\end{equation}
A stopped-target surrogate injects this direction into the generator.
The fake scorer receives a separate noise-prediction update on detached
samples; the generator update does not differentiate through the scorers.
The DMD reference also uses LPIPS-VGG regression to cached teacher images,
with weight 0.25. \emph{Both perceptual regression and this default weight
come from original DMD}; they are not our proposed contribution. Our LoRA
capacity and smaller teacher bank differ from its published full-model
system.

\paragraph{SwiftBrush.}
SwiftBrush trains a one-step generator with variational score distillation,
using a frozen teacher and an adapted fake
scorer~\cite{nguyen2024swiftbrush}. It learns from captions and generated
samples without a fixed teacher-image regression bank. This avoids the
paired-target preparation needed by DMD's regression term, while increasing
the importance of balancing generator and scorer optimization. We retain
its source-style optimizer, scorer and EMA settings but use a rank-32
SD1.5 generator and COCO captions. Its teacher guidance and precision are
reported separately rather than treated as matched to DMD.

\paragraph{InstaFlow.}
InstaFlow simplifies noise--image couplings through reflow before one-step
distillation~\cite{liu2024instaflow}. For noise $z$ and teacher endpoint $y$,
the path $x_\tau=(1-\tau)z+\tau y$ gives
\begin{equation}
L_{\mathrm{RF}}=\mathbb E\|v_\theta(x_\tau,\tau,c)-(y-z)\|_2^2.
\end{equation}
The one-step endpoint is $z+v_\theta(z,0,c)$. We convert the pretrained
epsilon prediction into velocity and use cached reflow couplings before
endpoint distillation. Reflow makes a simple inference path possible, but
coupling generation has a cost and errors in the couplings affect the
resulting student.

\begin{table*}[t]
\centering\small
\setlength{\tabcolsep}{4pt}
\begin{tabularx}{\textwidth}{lXXX}
\toprule
Pipeline & Main learning signal & Additional training machinery & Practical trade-off \\
\midrule
PD & Two-to-one transition targets & Repeated compression stages & Explicit targets; accumulated stage error \\
Consistency & Agreement along a trajectory & Teacher transition and EMA target & Direct endpoint map; boundary sensitivity \\
DMD & Distribution matching + LPIPS & Fake scorer and paired teacher bank & Image anchors; scorer and target costs \\
SwiftBrush & Variational score distillation & Adapted fake scorer & No paired-image bank; coupled optimization \\
InstaFlow & Reflow and endpoint regression & Cached noise--image couplings & Simple inference; coupling dependence \\
Ours & DMD + response direction alignment & Reuses paired DMD outputs & No extra inference call; added training objective \\
\bottomrule
\end{tabularx}
\caption{Learning signals and training requirements in the implemented pipelines.
All students use one NFE at inference. Equal inference cost does not imply
equal training cost or identical source-paper implementations.}
\label{tab:taxonomy}
\end{table*}

\subsection{Our direction-alignment extension: overview}
Our extension retains the DMD reference and adds an angular constraint
between student and teacher latent displacements. Two independent noises
share one caption. The student difference should point in the same direction
as the corresponding difference between cached teacher outputs, without
explicitly forcing the same magnitude. Existing paired-regression outputs
are reused, so the added term introduces no model parameters or additional
teacher/generator forward passes relative to that paired DMD update.
Inference remains unchanged.

Response matching itself has prior art, including geometric alignment of
local teacher and student responses~\cite{huang2026gad}. Our study tests a
particular direction-only formulation with independent cached pairs.
Appendix~\ref{sec:response} gives the objective, its relation to prior work,
fixed coefficients and the independent 36,000-update control. The main
report compares the selected model against the five standard pipelines.

\section{Experiments and analysis}
\subsection{Dataset and evaluation protocol}
We use COCO-2017~\cite{chen2015coco}, selecting the lowest-ID caption per
image. Two disjoint training pools of 50,000 and 68,287 captions cover the
118,287 training images; the later pool is revisited through shuffled
epochs. All 5,000 official validation images are held out from student
training. PD and consistency use real image--caption pairs. DMD and
SwiftBrush use captions with generated samples; DMD and InstaFlow also use
cached teacher targets.

Each quality evaluation generates one image per validation caption with
the same noise stream, seed 12345. All methods use 5,000 generated and
5,000 reference images with the same resizing and center-cropping protocol.
FID measures the discrepancy between real and generated Inception feature
statistics~\cite{heusel2017fid}; lower is better. Inception Score (IS)
measures class confidence and marginal coverage~\cite{salimans2016improved}.
CLIP ViT-B/32 measures image--text alignment as
$100\max(\cos(I,c),0)$, averaged over images~\cite{hessel2021clipscore}.
Higher IS and CLIP are better under these definitions. Published COCO scores
using other sampling or preprocessing protocols are not inserted into our
table.

All reported student checkpoints have exactly 36,000 cumulative generator
updates. Intermediate checkpoints and additional-update refinements do not
enter the quantitative tables or qualitative figures.

\subsection{Training budget and implementation settings}
The five standard students have 36,000 generator updates, effective batch
16, training seed 42 and EMA evaluation. DMD-family paired runs use
microbatch 2 with eight accumulation passes; SwiftBrush uses microbatch 1
with sixteen passes. The shared DMD teacher bank contains 5,000 outputs from
2,500 training captions and two noises per caption. Preparing it costs
$5{,}000\times64\times2=640{,}000$ teacher NFEs. Distribution matching uses
the full caption stream, while paired regression uses this smaller bank.

\begin{table}[t]
\centering\scriptsize
\setlength{\tabcolsep}{3pt}
\begin{tabular}{lcc}
\toprule
Setting & DMD family & SwiftBrush \\
\midrule
Generator rank / effective batch & 32 / 16 & 32 / 16 \\
Optimizer & AdamW & Adam \\
Generator learning rate & $10^{-4}$ & $10^{-6}$ \\
Fake-scorer learning rate & $10^{-4}$ & $10^{-3}$ \\
Generator schedule & Warmup/cosine & Constant \\
Warmup updates / weight decay & 100 / 0.01 & 0 / 0 \\
Gradient clipping & 1 & 1 \\
EMA decay & 0.995 & Up to 0.9999 \\
Scorer updates per G update & 1 & 1 \\
Fake-scorer rank & 32 & 64 \\
Teacher CFG / precision & 7.5 / BF16 & 4.5 / FP32 \\
Student inference CFG & 1 & 1 \\
\bottomrule
\end{tabular}
\caption{Settings for the 36,000-update standard pipelines. PD, consistency
and InstaFlow share the DMD-family generator settings without fake-scorer
updates. SwiftBrush retains its source-style scorer and EMA ramp. The
proposal's angular objective and teacher banks are described below.}
\label{tab:settings}
\end{table}

Our proposal starts from pretrained SD1.5 with fresh rank-32 student and
scorer adapters and receives 36,000 updates. Direction alignment is active
from the first update, with fixed coefficient $\lambda=0.298494$. It uses the DMD-family optimizer settings above
and 18,000 updates on each caption pool. Each phase has its own bank of
5,000 teacher outputs, giving 10,000 targets from 5,000 captions in total
and 1,280,000 teacher NFEs for preparation. The main table and all Ours
figures use this completed checkpoint. The appendix details the same
36,000-update experiment.

The comparison matches backbone, generator rank, baseline caption pools
and one-NFE student inference. It does not equalize total training compute:
PD has multiple stages, InstaFlow needs couplings, and DMD/SwiftBrush train
auxiliary scorers. SwiftBrush also uses different CFG, precision and
optimization (Table~\ref{tab:settings}). Ours and the original DMD reference
have equal generator-update counts but different initialization histories
and teacher banks. These distinctions matter when comparing complete recipes.

\subsection{Quantitative comparison}
\begin{table*}[t]
\centering\small
\input{generated/assignment2/concise_quality.tex}
\caption{\textbf{Full-validation quality including Ours.} Every score uses
all 5,000 COCO validation images. Bold marks the best student value for each
metric; the teacher is a reference with a larger inference budget. All
student generators have rank 32 and 36,000 updates. Ours is trained
independently with two phase-specific teacher banks; the original DMD
reference uses its earlier training setup.}
\label{tab:quality}
\end{table*}
Among the five standard pipelines, DMD has the lowest FID and highest IS and
CLIP (Table~\ref{tab:quality}). SwiftBrush is next in FID, followed by PD and
InstaFlow; the local consistency adaptation has the largest gap. DMD's
paired perceptual targets provide direct image-level anchors alongside its
distribution-matching signal. The regression ablation below examines that
component while holding the paired training setup fixed.

\paragraph{Our selected result.}
\input{generated/assignment2/selected_response_summary.tex}
Its IS is \ASelectedIS{}, compared with \AReferenceIS{} for DMD. The lower
observed FID comes with lower IS and text alignment. Both models have
36,000 generator updates, but their initialization histories and teacher
banks differ. \AIndependentControlStatus{}
The comparison with the original reference therefore does not isolate the
response term's effect.

The FID difference is \ASelectedFIDGain{} points, or
\ASelectedFIDGainPercent{}\% of the reference value. IS and CLIP change by
\ASelectedISChange{} and \ASelectedCLIPChange{}, respectively. Thus the
advantage is specific to FID. FID compares marginal image-feature
statistics, while CLIP checks image--caption compatibility; an improvement
in the former need not improve the latter. The added objective constrains
responses to noise under a fixed caption and does not directly supervise
caption relations or object counts. Appendix~\ref{sec:response} analyzes
the corresponding gradient and its interaction with the DMD objective.

\paragraph{What the baseline ranking explains.}
SwiftBrush avoids a paired-image bank but must balance generator and fake
scorer updates under its own guidance and learning rates. PD spreads the
budget across several compression stages, making the quality of its
intermediate targets relevant. InstaFlow depends on cached couplings and
our conversion from epsilon prediction to velocity. Consistency is
sensitive to its boundary and loss choices; the local VP/latent-L2 variant
is not the original EDM/perceptual system. The scores identify strengths
and weaknesses of these local recipes without establishing a universal
ranking of the published methods.

The multi-step teacher remains better in FID and CLIP than every student.
One-step sampling removes repeated denoiser calls but does not imply a
128-fold end-to-end speedup: text encoding, VAE decoding, batching and
hardware utilization still contribute. Concurrent workloads make the
recorded wall times unsuitable for a controlled latency ranking.

\subsection{DMD regression ablation and setting analysis}
\begin{table}[t]
\centering\small
\input{generated/assignment2/regression_control.tex}
\caption{Matched DMD regression-space ablation at 36,000 total updates.
Both variants use the same paired bank, initialization, optimizer state and
caption stream. Regression weight is fixed at 0.25.}
\label{tab:regression}
\end{table}
We replace DMD's LPIPS regression with latent MSE while retaining its
score-based objective. The two variants share their first 18,000 updates
and then receive 18,000 updates with the specified regression loss.
Both use two cached noises per caption and the same pair grouping.
LPIPS reduces FID from \AControlFID{} to \APerceptualFID{}, a difference of
\ARegressionGain{}, and improves IS and CLIP. This supports decoded-image
perceptual supervision over latent coordinate error in the tested recipe.

The coefficient 0.25 is inherited from DMD rather than selected as a novel
hyperparameter. Holding its scalar value fixed does not equalize the
magnitudes of LPIPS and MSE gradients, so the experiment compares the
resulting regression recipes rather than isolating feature representation
alone. Likewise, the optimizer and CFG differences in
Table~\ref{tab:settings} are reported settings, not individually established
causes of the cross-method ranking.

For our extension, the angular coefficient remains fixed throughout the
36,000-update run. The reported experiments do not establish its optimal
value. A matched zero-response run uses the same fresh adapters, teacher
banks and training schedule; its final scores are included only after
full validation finishes. This control addresses the added response term
separately from the inherited perceptual regression.

\subsection{Native samples and text control}
Figure~\ref{fig:main} compares identical prompts and noise across all six
students and the teacher. The \textbf{Ours} column contains outputs from the
same selected model used in Table~\ref{tab:quality}. DMD and Ours preserve
recognizable subjects and coherent layouts in these examples; their
local texture and framing differ. Other local pipelines show varying
levels of blur or structural error. These selected examples illustrate
appearance but do not establish a general visual winner between DMD and
Ours.

The teddy-bear prompt also shows a shared limitation: recognizable objects
do not guarantee exact counts, colors or attributes. All 5,000 generated
validation images enter the quality measurements, including outputs that
are less successful than the displayed examples. This prevents the selected
figure from determining the numerical ranking.

\begin{figure*}[p]
\centering
\includegraphics[width=0.76\textwidth]{generated/assignment2/prompt_control.pdf}
\caption{\textbf{Text control with Ours}, the selected
\ASelectedUpdates{}-update model, at one NFE. Prompts follow ``A DSLR
photograph of a \{animal\} \{action\}, sitting at an outdoor cafe table.''
The grid varies four animals and three actions with the same initial noise.
\AControlSelectionNote{}}
\label{fig:control}
\end{figure*}

Figure~\ref{fig:control} varies animal identity and action under fixed
noise. The outputs demonstrate the selected model's ability to change
subject appearance and activity while retaining a related scene layout.
The grid is a qualitative conditioning probe; it is not a measured
compositional-success benchmark. Small objects and exact interactions may
remain ambiguous even when the intended animal is recognizable.

\begin{figure*}[p]
\centering
\includegraphics[width=0.94\textwidth]{generated/assignment2/prompt_interpolation.pdf}
\caption{\textbf{Prompt-embedding interpolation with Ours.} Dog-to-cat weights
are $0,0.25,0.5,0.75,1$. Endpoint prompts describe an animal reading at an
outdoor cafe. Each row fixes the noise and uses the selected
\ASelectedUpdates{}-update model at one NFE. \AInterpolationSelectionNote{}}
\label{fig:interpolation}
\end{figure*}

Figure~\ref{fig:interpolation} interpolates the CLIP hidden states as
$e_\alpha=(1-\alpha)e_{\mathrm{dog}}+\alpha e_{\mathrm{cat}}$ and passes the
result through the generator. Noise stays unchanged within each row.
The images are native outputs, not pixel blends. The endpoints change the
animal identity, while intermediate conditions can mix features. This
illustrates local behavior in conditioning space without establishing
linear semantic geometry or reliable attribute control in general.

\clearpage
\subsection{Limitations and future extensions}
\paragraph{Capacity and compute.}
All students receive only 36,000 updates with rank-32 adapters.
SwiftBrush reports 65,000 updates, while consistency-model experiments
include 600,000--1,000,000 iterations across datasets and
variants~\cite{nguyen2024swiftbrush,song2023consistency}. Different
architectures, batch sizes and teacher costs prevent direct compute
equivalence. Our results characterize the project budget, not the methods'
performance ceilings.

\paragraph{Selection and generalization.}
We use one training seed and one quantitative dataset. Validation images
are held out from student training but reused for development and selection.
Teacher-pretraining overlap with COCO is unknown. The original DMD bank
covers 2,500 captions; the proposal's two banks cover 5,000. Repeated
training seeds, the independent zero-response control, an untouched test
set and a second caption distribution would test the proposal's robustness.

\paragraph{Evaluation scope.}
FID, IS and CLIP do not directly measure object counting, spatial relations
or attribute binding.
Object-level evaluation and blinded human comparisons would complement
these metrics. A dedicated efficiency study should isolate GPU workloads
and measure full inference latency, peak memory and total training cost,
including teacher-target preparation.

\section{Conclusion}
Among five standard pipelines under a shared rank-32 SD1.5 architecture,
DMD with its original LPIPS regression is strongest; a matched ablation
supports that component. Our direction-alignment extension reaches FID
\ASelectedFID{} versus \ADMDFID{} for the original DMD reference at 36,000
generator updates each, with lower IS and CLIP. Native samples and conditioning
probes complement these measurements. The appendix details the design and
36,000-update training protocol. The independent control and repeated seeds
are needed to assess how much of the FID difference comes from the added objective.

{\small
\bibliographystyle{ieeenat_fullname}
\bibliography{assignment2_references}
}
\clearpage
% Count the main report after its floats and references have been shipped out.
\typeout{ASSIGNMENT2-MAIN-PAGES:\the\numexpr\value{page}-1\relax}
\appendix
\section{Our proposed direction-alignment method}
\label{sec:response}
\begin{figure*}[t]
\centering
\includegraphics[width=\textwidth]{generated/assignment2/response_alignment_method.pdf}
\caption{\textbf{Our proposed direction-alignment design.} Two independent noises share a
caption and cached teacher targets. The added angular term aligns student
and teacher latent-displacement directions; it does not match their
magnitudes. Existing DMD and LPIPS objectives remain. Both student outputs
receive gradients, teacher targets are fixed, and inference uses one denoiser
call. The vectors are schematic illustrations of the objective.}
\label{fig:method}
\end{figure*}
\begin{samepage}
Pointwise regression does not explicitly constrain how outputs change
between noise inputs. For independent cached noises $z_1,z_2$ with the same
caption, define the finite latent displacements
\begin{align}
d_S &= G_\theta(z_2,c)-G_\theta(z_1,c),\\
d_T &= T(z_2,c)-T(z_1,c),\\
L_{\mathrm{angle}} &= 1-\frac{\langle d_S,d_T\rangle}
 {\|d_S\|_2\|d_T\|_2},\\
L &= L_{\mathrm{DMD}}+0.25L_{\mathrm{LPIPS}}+
\lambda L_{\mathrm{angle}}.
\end{align}
\end{samepage}
Teacher targets are detached; numerical guards handle vanishing norms.
The added term corrects direction without explicitly forcing the teacher's
response magnitude. For nonzero displacements, its two latent gradients
sum to zero and are perpendicular to $d_S$. These are first-order properties
of this term, not guarantees about the combined Adam update.

Response alignment has prior work: GAD matches local teacher and student
responses through squared vector error on nearby perturbed
noises~\cite{huang2026gad}. Our experiment tests angular alignment using
independent cached pairs. These finite displacements are not local Jacobian
estimates. The contribution under study is the usefulness of this particular
direction-only choice, not the invention of response matching or cosine loss.
Paired-regression outputs are reused, adding no model parameters or extra
teacher/generator forward passes relative to the paired DMD reference.


\subsection{Independent 36,000-update training}
\input{generated/assignment2/selected_response_settings.tex}
The angular coefficient is held constant throughout training, rather than
adjusted from validation scores during the run. We do not claim that it is
optimal among 36,000-update configurations. The LPIPS weight 0.25 is inherited
from DMD and is shared with the zero-response control.

Each phase uses 2,500 training captions and two cached noises per caption.
Distribution matching samples the full phase caption pool, while paired
regression and direction alignment use the cached pairs. Teacher targets
are generated with 64 DDIM steps and CFG 7.5. Preparing both phase banks
therefore uses 10,000 images and 1,280,000 teacher NFEs. The two phases
contribute 18,000 student updates each; the final checkpoint has exactly
36,000 cumulative student updates and 36,000 scorer updates.

The student uses one conditional U-Net evaluation at inference, followed
by the frozen VAE decoder. Neither the fake scorer nor the paired teacher
bank is needed for generation. The main comparison, animal/action grid
and embedding interpolation all use this same terminal checkpoint.
\input{generated/assignment2/selected_response_decision.tex}

\subsection{Geometry and interpretation}
Write $u=d_S/\|d_S\|$ and $v=d_T/\|d_T\|$. Away from the numerical guard,
\begin{equation}
\nabla_{d_S}L_{\rm angle}
=-\frac{(I-uu^\top)v}{\|d_S\|}.
\end{equation}
The projection removes the radial component, so
$d_S^\top\nabla_{d_S}L_{\rm angle}=0$. Since the two endpoints enter with
opposite signs, their angular gradients sum to zero. This describes the
loss's gradient in latent coordinates. Shared network parameters, the
other objectives and Adam's coordinate scaling mean the actual parameter
update need not preserve either the displacement norm or the midpoint.

To compare direction alignment with a full displacement error, set
$r=\|d_S\|/\|d_T\|$. The relative squared error is
\begin{equation}
\frac{\|d_S-d_T\|^2}{\|d_T\|^2}
=(r-1)^2+2r(1-u^\top v).
\end{equation}
It couples a radial penalty to an angular penalty weighted by the current
norm ratio. Direction alignment instead gives orientation its own fixed
coefficient. This algebra explains the design choice; it is not a measured
comparison against a separately trained vector-loss model at 36,000 updates.

The angular gradient has norm $|\sin\phi|/\|d_S\|$ above the numerical guard,
where $\phi$ is the displacement angle. A small student response can amplify
this signal. The implementation computes differences in FP32, excludes
pairs with $\|d_T\|\leq10^{-6}$, clamps the student norm below at $10^{-6}$,
and applies the shared generator gradient clipping. The gradient is zero
at parallel and exactly antiparallel directions; angular alignment alone
does not guarantee a successful optimization trajectory.

Direction alignment is invariant to a positive scalar rescaling of $d_S$.
It therefore does not determine the correct displacement magnitude, image
appearance or coverage on its own. LPIPS anchors decoded images, and the
score objective supplies a distribution-level signal. The experiment tests
the combined recipe; a low angular loss is not a standalone quality metric.
Independent paired noises define finite displacements, so this objective
does not estimate a local Jacobian or imply global geometric agreement.

\subsection{Interaction with the DMD learning signals}
LPIPS constrains each decoded endpoint, and the angular term directly
compares the relative latent orientation of the pair. These constraints
overlap but are expressed in different spaces: a decoded-image perceptual
distance and a flattened-latent direction. The teacher difference is fixed;
only the two student branches receive the response gradient. The fake
scorer is still trained by its denoising objective and receives no response
gradient. The same cached pairs supply the regression and response terms.

This reuse keeps the number of teacher and student forward passes unchanged
relative to the paired DMD update. It adds latent-vector reductions and
their gradients, with no added model parameters or inference calls. We
do not report a measured training-speed gain: equal forward-pass counts
do not establish equal wall-clock time. A one-NFE student still requires
text encoding and VAE decoding in the complete inference pipeline.

The animal/action grid and text interpolation probe conditioning behavior,
whereas the response objective changes noise at a fixed caption. They are
therefore complementary illustrations, rather than direct evidence that
the angular constraint improves text control. A matched zero-response
experiment and repeated seeds are needed to separate the added loss from
the complete recipe and measure the reliability of its contribution.

\input{generated/assignment2/independent_results.tex}
\input{generated/assignment2/recent_baselines.tex}

\end{document}
